\chapter{Herencia y Polimorfismo}

\section{Herencia}
La \textit{herencia} es un mecanismo fundamental de la POO que permite crear nuevas clases basadas en clases existentes.
La clase original se denomina \textit{superclase} o \textit{clase padre}, mientras que la nueva clase se llama
\textit{subclase} o \textit{clase hija}.

La subclase hereda automáticamente todos los atributos y métodos de la superclase, lo que fomenta la reutilización de código
y permite establecer jerarquías conceptuales del tipo \textit{``es-un''} (por ejemplo, un colectivo \textit{es un} vehículo).

\begin{figure}[!ht]
  \centering
  \begin{tikzpicture}[node distance=1.5cm, align=center]
    \node [draw, rectangle, fill=blue!20, minimum width=3.5cm, minimum height=1.2cm] (padre) 
      {\textbf{Superclase}\\ \texttt{Vehicle}};
    
    \node [draw, rectangle, fill=green!20, minimum width=2.5cm, minimum height=1cm, below left=1cm and 0.5cm of padre] (hijo1) 
      {\texttt{Car}};
    
    \node [draw, rectangle, fill=green!20, minimum width=2.5cm, minimum height=1cm, below=1cm of padre] (hijo2) 
      {\texttt{Bus}};
    
    \node [draw, rectangle, fill=green!20, minimum width=2.5cm, minimum height=1cm, below right=1cm and 0.5cm of padre] (hijo3) 
      {\texttt{Truck}};
    
    \draw[<-, thick] (padre) -- (hijo1);
    \draw[<-, thick] (padre) -- (hijo2);
    \draw[<-, thick] (padre) -- (hijo3);
  \end{tikzpicture}
  \caption{Jerarquía de herencia a partir de la superclase Vehicle.}
  \label{fig:herencia_vehiculos}
\end{figure}

\begin{lstlisting}[caption={Sintaxis básica de herencia en Python.}, label={list:herencia_basica}]
class Vehicle:
    def __init__(self, name, max_speed):
        self.name = name
        self.max_speed = max_speed

    def display_info(self):
        print(f"Vehiculo: {self.name}, Velocidad Maxima: {self.max_speed} km/h")

# Bus hereda de Vehicle
class Bus(Vehicle):
    pass

b = Bus("Mercedes Benz", 120)
b.display_info()  # Metodo heredado de Vehicle
\end{lstlisting}

\section{Sobrescritura de Métodos y Uso de \texttt{super()}}
En muchas situaciones, una subclase necesita adaptar o extender el comportamiento de un método heredado.
Este proceso se denomina \textit{sobrescritura de métodos} (\textit{method overriding}).

Para reutilizar la lógica de la superclase en lugar de reescribirla por completo, Python proporciona la función \texttt{super()}.

\begin{lstlisting}[caption={Uso de \texttt{super()} para extender el constructor y métodos.}, label={list:super_ejemplo}]
class Vehicle:
    def __init__(self, name, max_speed, mileage):
        self.name = name
        self.max_speed = max_speed
        self.mileage = mileage

    def fare(self):
        return self.mileage * 10

class Bus(Vehicle):
    def __init__(self, name, max_speed, mileage, capacity=50):
        # Llama al constructor de la superclase Vehicle
        super().__init__(name, max_speed, mileage)
        self.capacity = capacity

    def fare(self):
        # Sobrescribe el metodo agregando un 10% de mantenimiento
        base_fare = super().fare()
        total_fare = base_fare + (base_fare * 0.10)
        return total_fare

bus = Bus("Escolar", 90, 100, 40)
print(f"Tarifa del colectivo: ${bus.fare()}")
\end{lstlisting}

\section{Inspección de Tipos y Jerarquías}
Python dispone de funciones integradas para consultar la pertenencia de un objeto a una clase o jerarquía:

\begin{itemize}[leftmargin=*]
  \item \texttt{type(obj)}: retorna la clase exacta a la cual pertenece un objeto.
  \item \texttt{isinstance(obj, Clase)}: comprueba si un objeto es una instancia de la clase especificada o de una subclase derivada de ella.
  \item \texttt{issubclass(Subclase, Superclase)}: verifica si una clase hereda de otra superclase.
\end{itemize}

\begin{lstlisting}[caption={Uso de \texttt{isinstance()} e \texttt{issubclass()}.}, label={list:inspeccion_tipos}]
b = Bus("Urbano", 80, 50)

print(type(b))                       # <class '__main__.Bus'>
print(isinstance(b, Bus))           # True
print(isinstance(b, Vehicle))       # True (por herencia)
print(issubclass(Bus, Vehicle))     # True
\end{lstlisting}

\section{Polimorfismo}
El \textit{polimorfismo} es la capacidad que poseen distintos objetos de responder al mismo mensaje o llamada
de método, ejecutando cada uno su propia implementación específica.

En Python, el polimorfismo se aplica frecuentemente mediante el principio de \textit{Duck Typing} (``Si camina como
un pato y suena como un pato, entonces es un pato''). Esto significa que a una función no le importa la clase exacta
del objeto recibido, sino que este responda a los métodos invocados.

\begin{lstlisting}[caption={Polimorfismo con jerarquía de animales.}, label={list:polimorfismo_animales}]
class Dog:
    def speak(self):
        return "Guau!"

class Cat:
    def speak(self):
        return "Miau!"

def hacer_sonar(animal):
    # Invoca speak() de forma polimorfica sin importar el tipo exacto
    print(animal.speak())

mascotas = [Dog(), Cat()]
for m in mascotas:
    hacer_sonar(m)
\end{lstlisting}

Un ejemplo aplicado al cálculo de salarios en un sistema de recursos humanos ilustra el valor del polimorfismo:

\begin{lstlisting}[caption={Polimorfismo aplicado al cálculo de salarios.}, label={list:polimorfismo_empleados}]
class Employee:
    def __init__(self, name):
        self.name = name

    def calculate_salary(self):
        raise NotImplementedError("Las subclases deben implementar este metodo.")

class FullTimeEmployee(Employee):
    def __init__(self, name, monthly_salary):
        super().__init__(name)
        self.monthly_salary = monthly_salary

    def calculate_salary(self):
        return self.monthly_salary

class PartTimeEmployee(Employee):
    def __init__(self, name, hours_worked, hourly_rate):
        super().__init__(name)
        self.hours_worked = hours_worked
        self.hourly_rate = hourly_rate

    def calculate_salary(self):
        return self.hours_worked * self.hourly_rate

empleados = [
    FullTimeEmployee("Carlos", 350000),
    PartTimeEmployee("Ana", 80, 2500)
]

for e in empleados:
    print(f"Empleado: {e.name}, Salario: ${e.calculate_salary()}")
\end{lstlisting}
